\begin{center}
[4].\\
\textsc{Gauss} an \textsc{Olbers}.\quad Göttingen, 13.~Januar 1821.
\end{center}

\par\noindent\hspace*{\parindent}.\ .\ .\ .\ .\ .\quad Sollte es in Zukunft bei der wirklichen Messung mit allen äußeren Umständen besser gehen, als es bisher den Anschein hat, so glaube ich, dass ich wohl Freude an der Arbeit haben könnte, und dann würde ich mich auch recht gern einer Erweiterung der Triangulation nach Westen, falls sie mir aufgetragen würde, unterziehen.
Die Anschließung an die \textsc{Krayenhoff}schen Dreiecke ist allerdings wünschenswert, allein wo sind denn diese zu finden?
Ich weiß nicht, ob sie irgendwo gedruckt sind, und der schlechte Erfolg mit den \textsc{Epailly}schen Dreiecken macht mich ganz mutlos.
Auch \textsc{Laplace}, an den ich
\par\clearpage
\noindent368\hfill{\small\textbf{BRIEFWECHSEL MIT OLBERS.}}\par\medskip
\noindent vor etwa 9 Wochen geschrieben habe, hat mir gar nicht geantwortet.
Meiner Meinung nach sollten alle gut gemessenen Dreiecke 1.~Ordnung als etwas betrachtet werden, worauf das ganze Publikum Anspruch hat, und nach und nach sollte ganz Europa mit solchen Dreiecken überzogen werden.
Ich habe mir schon seit Jahren eine eigene Methode entworfen, wie solche Messungen am zweckmäßigsten behandelt werden können; denn alles, was ich darüber gelesen habe, finde ich herzlich wertlos.
So haben sich z.\ B.\ viele Mathematiker große Mühe mit der Aufgabe gegeben, aus Abständen vom Meridian und Perpendikel die Länge und Breite zu berechnen, mit Rücksicht auf die elliptische Gestalt der Erde, während, so viel ich weiß, niemand vorher gefragt hat:

\par\noindent\hspace*{\parindent}1.\ wie denn jene Abstände, so verstanden, wie man sie gewöhnlich versteht, aus der Messung mit ebenso großer Schärfe gefunden werden können; denn es scheint, dass die meisten diese Rechnung wie in der Ebene führen, oder doch ganz unrichtige oder unbrauchbare Vorschriften dafür geben;

\par\noindent\hspace*{\parindent}2.\ ob es denn überhaupt nur zweckmäßig sei, die \textls[120]{so verstandenen Abstände} zu gebrauchen, da es entschieden ist, dass, wenn man sie hinlänglich scharf aus den Dreiecken ableiten will, dies nur durch höchst beschwerliche Rechnungen geschehen kann, so wie man aus ihnen nur mit vieler Mühe wieder zu den Längen und Breiten herabsteigt.
Das Ganze würde nur ein »die Pferde hinter den Wagen spannen« sein.
»Soll etwas Brauchbares \textls[120]{zwischen} die Dreiecke und die Längen und Breiten gesetzt werden, so muss es etwas ganz anderes wie jene, so wie gewöhnlich verstanden, Koordinaten sein.«
Wie dies bei meiner Theorie geschieht, kann ich hier freilich nicht umständlich ausführen; nur so viel bemerke ich, dass das, was ich zwischen die Dreiecke und die Längen und Breiten setze, \textls[120]{diejenigen Koordinaten sind, 1)~mit denen am zweckmäßigsten jeder Punkt in einer Ebene dargestellt werden kann}.
Diese Koordinaten folgen höchst bequem und leicht aus den gemessenen Dreiecken, und ohne eine sehr genaue Kenntnis der Abplattung der Erde vorauszusetzen, und 2)~aus ihnen folgt wieder ebenso leicht die Länge und Breite, natürlich indem man die Abplattung kennen muss.
Ich habe die Absicht, diese Theorie, wo nicht früher, doch mit meinen künftigen Messungen bekannt zu machen und bitte vorerst, diese angedeuteten Ideen noch für sich zu behalten.
Sehr gern würde ich sie nicht bloß auf die hannoverschen Dreiecke, sondern auf alle anderen damit in Verbindung kommenden anwenden
\par\clearpage
\noindent{\small\textbf{ZUR HANNOVERSCHEN TRIANGULATION.}}\hfill369\par\medskip
\noindent und so eine \textls[120]{Description géométrique} eines großen Teils von Europa geben, wenn ich durch Mitteilung gehörig unterstützt würde.
Aber!!

\par\noindent\hspace*{\parindent}Herr v.\ \textsc{Müffling} hat mir doch seine 15 Dreiecke vom Rhein bis Seeberg mitgeteilt.
Vorläufig, aber freilich nur sehr roh, habe ich bereits Göttingen angeschlossen.
Nämlich 1812 habe ich auf dem Hanstein, dessen Lage gegen Göttingen näherungsweise aus meinen Winkelmessungen in hiesiger Gegend folgt, die Winkel zwischen Göttingen, Brocken und der Boineburg (2.~\textsc{Müffling}scher Punkt) gemessen, freilich auf mehrere Minuten ungewiss; doch glaube ich, dass die Eintragung Göttingens, die hieraus folgt, wohl auf 100 Meter beinahe zuverlässig ist.
Es folgt daraus: Längenunterschied zwischen der Göttinger und Seeberger Sternwarte in Zeit \mbox{3\textsuperscript{m}8\rlap{\textsuperscript{s}},7}, was sehr nahe mit den astronomischen Bestimmungen zutrifft.
Paris wäre hiernach, wenn es \mbox{33\textsuperscript{m}35\textsuperscript{s}} westlich von Seeberg liegt, \mbox{30\textsuperscript{m}26\rlap{\textsuperscript{s}},3} westlich von Göttingen.
Die neue Sternwarte liegt \mbox{1\rlap{\textsuperscript{s}},9} östlich von der alten, die ich früher immer \mbox{30\textsuperscript{m}23\(\tfrac{1}{2}\)\textsuperscript{s}} von Paris setzte.
